\documentclass[11pt]{article}

\usepackage[margin=0.5cm]{geometry}
\usepackage{amsmath}
\usepackage{amssymb}
\usepackage{enumitem}
\usepackage{array}
\usepackage{graphicx}
\usepackage{svg}

\begin{document}

\twocolumn

\title{Problem Set 3, COSC330/PHYS306\\
Chapters 3 and 4}
\author{Prof. Jordan C. Hanson}
\maketitle

%===========================================================
% Chapter 3
%===========================================================
\section{Chapter 3}

\begin{enumerate}

\item \textbf{[Section 3-9, Exercise 39]}
Examine the conditions indicated in Figure 3--92, and identify the faulty gates.

\begin{center}
\begin{tabular}{cc}
\includesvg[width=0.43\linewidth]{fig_3_92a} &
\includesvg[width=0.43\linewidth]{fig_3_92b} \\
\textbf{(a)} & \textbf{(b)} \\[0.15cm]
\includesvg[width=0.43\linewidth]{fig_3_92c} &
\includesvg[width=0.43\linewidth]{fig_3_92d} \\
\textbf{(c)} & \textbf{(d)} \\[0.15cm]
\includesvg[width=0.43\linewidth]{fig_3_92e} &
\includesvg[width=0.43\linewidth]{fig_3_92f} \\
\textbf{(e)} & \textbf{(f)}
\end{tabular}

\vspace{0.08cm}
\textbf{FIGURE 3--92}
\end{center}

\item \textbf{[Section 3-9, Exercise 40]}
Determine the faulty gates in Figure 3--93 by analyzing the timing diagrams.

\begin{center}
\includesvg[width=0.58\linewidth]{fig_3_93a}
\hfill
\includesvg[width=0.36\linewidth]{fig_3_93a_gate}

\textbf{(a)}

\vspace{0.12cm}
\includesvg[width=0.58\linewidth]{fig_3_93b}
\hfill
\includesvg[width=0.36\linewidth]{fig_3_93b_gate}

\textbf{(b)}

\vspace{0.12cm}
\includesvg[width=0.58\linewidth]{fig_3_93c}
\hfill
\includesvg[width=0.36\linewidth]{fig_3_93c_gate}

\textbf{(c)}

\vspace{0.12cm}
\includesvg[width=0.58\linewidth]{fig_3_93d}
\hfill
\includesvg[width=0.36\linewidth]{fig_3_93d_gate}

\textbf{(d)}

\vspace{0.08cm}
\textbf{FIGURE 3--93}
\end{center}

\end{enumerate}

%===========================================================
% Chapter 4
%===========================================================
\section{Chapter 4}

\begin{enumerate}
\setcounter{enumi}{2}

\item \textbf{[Section 4-3, Exercise 9]}
Apply DeMorgan's theorems to each expression:
\begin{enumerate}[label=(\alph*)]
    \item $\overline{A+B}$
    \item $\overline{\overline{A}B}$
    \item $\overline{A+B+C}$
    \item $\overline{ABC}$
    \item $\overline{A(B+C)}$
    \item $\overline{AB+CD}$
    \item $\overline{\overline{A}B+CD}$
    \item $\overline{(A+\overline{B})(\overline{C}+D)}$
\end{enumerate}

\item \textbf{[Section 4-4, Exercise 13]}
Write the Boolean expression for each of the logic circuits in Figure 4--57.

\begin{center}
\begin{tabular}{cc}
\includesvg[width=0.43\linewidth]{fig_4_57a} &
\includesvg[width=0.43\linewidth]{fig_4_57b} \\
\textbf{(a)} & \textbf{(b)} \\[0.18cm]
\includesvg[width=0.43\linewidth]{fig_4_57c} &
\includesvg[width=0.43\linewidth]{fig_4_57d} \\
\textbf{(c)} & \textbf{(d)}
\end{tabular}

\vspace{0.08cm}
\textbf{FIGURE 4--57}
\end{center}

\item \textbf{[Section 4-4, Exercise 15]}
Draw the logic circuit represented by each expression:
\begin{enumerate}[label=(\alph*)]
    \item $A(B+C)$
    \item $AB+ABC$
    \item $\overline{ABC}+D$
    \item $A+B(C+DB+\overline{C})$
\end{enumerate}

\item \textbf{[Section 4-4, Exercise 18]}
Construct a truth table for each of the following Boolean expressions:
\begin{enumerate}[label=(\alph*)]
    \item $A+B$
    \item $AB$
    \item $A+B(C+D)$
    \item $A+B+C$
    \item $AB+BC$
\end{enumerate}

\item \textbf{[Section 4-4, Exercise 22]}
Determine which of the logic circuits in Figure 4--59 are equivalent.

\begin{center}
\begin{tabular}{cc}
\includesvg[width=0.43\linewidth]{fig_4_59a} &
\includesvg[width=0.43\linewidth]{fig_4_59b} \\
\textbf{(a)} & \textbf{(b)} \\[0.18cm]
\includesvg[width=0.43\linewidth]{fig_4_59c} &
\includesvg[width=0.43\linewidth]{fig_4_59d} \\
\textbf{(c)} & \textbf{(d)}
\end{tabular}

\vspace{0.08cm}
\textbf{FIGURE 4--59}
\end{center}

\item \textbf{[Section 4-7, Exercise 31]}
Develop a truth table for each of the following standard SOP expressions:
\begin{enumerate}[label=(\alph*)]
    \item $ABC+\overline{A}BC+AB\overline{C}$
    \item $XYZ+X\overline{Y}Z+\overline{X}YZ+\overline{X}Y\overline{Z}
          +X\overline{Y}\overline{Z}$
\end{enumerate}

\item \textbf{[Section 4-7, Exercise 36]}
For each truth table in Table 4--15, derive a standard SOP and a standard POS expression.

\begin{center}
\scriptsize
\renewcommand{\arraystretch}{1.10}

\begin{tabular}{c c c|c}
\multicolumn{3}{c|}{$ABC$} & $X$\\ \hline
0&0&0&0\\0&0&1&1\\0&1&0&0\\0&1&1&0\\
1&0&0&1\\1&0&1&1\\1&1&0&0\\1&1&1&1
\end{tabular}
\hspace{0.25cm}
\begin{tabular}{c c c|c}
\multicolumn{3}{c|}{$ABC$} & $X$\\ \hline
0&0&0&0\\0&0&1&0\\0&1&0&0\\0&1&1&0\\
1&0&0&0\\1&0&1&1\\1&1&0&1\\1&1&1&1
\end{tabular}

\vspace{0.05cm}
\textbf{(a)}\hspace{2.15cm}\textbf{(b)}

\vspace{0.18cm}

\begin{tabular}{c c c c|c}
\multicolumn{4}{c|}{$ABCD$} & $X$\\ \hline
0&0&0&0&1\\0&0&0&1&1\\0&0&1&0&0\\0&0&1&1&1\\
0&1&0&0&0\\0&1&0&1&1\\0&1&1&0&1\\0&1&1&1&0\\
1&0&0&0&0\\1&0&0&1&1\\1&0&1&0&0\\1&0&1&1&0\\
1&1&0&0&1\\1&1&0&1&0\\1&1&1&0&0\\1&1&1&1&0
\end{tabular}
\hspace{0.12cm}
\begin{tabular}{c c c c|c}
\multicolumn{4}{c|}{$ABCD$} & $X$\\ \hline
0&0&0&0&0\\0&0&0&1&0\\0&0&1&0&1\\0&0&1&1&0\\
0&1&0&0&1\\0&1&0&1&1\\0&1&1&0&0\\0&1&1&1&1\\
1&0&0&0&0\\1&0&0&1&0\\1&0&1&0&0\\1&0&1&1&1\\
1&1&0&0&1\\1&1&0&1&0\\1&1&1&0&0\\1&1&1&1&1
\end{tabular}

\vspace{0.05cm}
\textbf{(c)}\hspace{2.30cm}\textbf{(d)}

\vspace{0.06cm}
\textbf{TABLE 4--15}
\end{center}

\item \textbf{[Section 4-9, Exercise 40]}
Use a Karnaugh map to find the minimum SOP form for each expression:
\begin{enumerate}[label=(\alph*)]
    \item $\overline{A}\,\overline{B}C
          +A\overline{B}C
          +\overline{A}B\overline{C}$
    \item $AC(\overline{B}+C)$
    \item $\overline{A}(BC+B\overline{C})
          +A(BC+B\overline{C})$
    \item $\overline{A}\,\overline{B}C
          +A\overline{B}\,\overline{C}
          +\overline{A}B\overline{C}
          +AB\overline{C}$
\end{enumerate}

\end{enumerate}

\end{document}
